\section[Stage 001]{Stage 001: Geometry Lift and Linearized PDE Skeleton}
\label{stage:001}

\stagefield{Part / anchor}{Part I; primary anchors \resultanchor{MTDC-T1}, \resultanchor{MTDC-T2}, and \resultanchor{MTDC-T4}.}
\claimstatus{The parent-field import is \StatusExact{} inside the declared parent model.  The shape-field lift and minimal quadratic wall action are \StatusEffective{}.  The modal split and the linearized equations are \StatusExactClosure{} once that lift is adopted.}

\stagefield{Purpose}{Stage~001 replaces the old finite-dimensional geometry variables \(a(t)\) and \(L(t)\) by a distributed throat geometry variable, while preserving those variables as collective moments.  The stage also introduces the first linearized matter/gauge/geometry system and the harmonic decomposition in which the scalar lane, grouped real \(P_2\) lane, geometry lane, and higher modes become explicit.}

\stagefield{Inputs}{The stage imports the exact \((4+1)\)-dimensional parent theory: gauged GNLS matter, localized Maxwell field, corrected charge ontology, and confinement potential.  It also imports the lower-order requirement that the moving-throat PDE expose the already-visible \(P_0\oplus P_2\) support sector, the separate geometry-completion channel, and the mixed Maxwell channels retained outside the strict zero-mode brane limit.}

\stagefield{Verification}{SymPy audit: \StageFile{scripts/moving_throat_pde_stage001_geometry_lift_sympy_audit.py}.  Mathematica audit: \StageFile{mathematica/moving_throat_pde_stage001_geometry_lift_mathematica_audit.wl}.  The audits check the monopole/\(P_2\) harmonic bookkeeping, the confinement chain-rule sign, and the densitized-versus-weighted wall-action convention used by the ledger.}

\subsection*{Derivation}

\paragraph{1. Shape-field lift.}
Let
\begin{equation}
\label{eq:app-stage001-sigma}
\Sigma(\mathbf X,t)=r-R(\Omega,w,t),
\qquad
r=\sqrt{x^2+y^2+z^2},
\qquad
\Omega=\mathbf x/r\in S^2.
\end{equation}
The finite throat surface is the level set
\begin{equation}
\label{eq:app-stage001-surface}
\Sigma(\mathbf X,t)=0.
\end{equation}
The working sign convention is
\begin{equation}
\Sigma>0\quad\text{outside the throat},
\qquad
\Sigma<0\quad\text{inside the support/interior region}.
\end{equation}
This choice keeps the level-set coupling to the bulk fields while making the mouth-sphere harmonics explicit.  It is the minimum lift needed to avoid treating the grouped real \(P_2\) coefficients as purely symbolic residuals.

\paragraph{2. Reference stationary branch.}
Take a stationary reference surface
\begin{equation}
\label{eq:app-stage001-reference}
\Sigma_0(\mathbf X)=r-R_0(w)=0,
\end{equation}
with mouth at \(w=0\), finite depth \(0\le w\le L_0\), mouth radius
\begin{equation}
\label{eq:app-stage001-a0}
a_0=R_0(0),
\end{equation}
and a bottom closure such as \(R_0(L_0)=0\) or an equivalent regularity condition.  The old collective variables are recovered as lowest moments:
\begin{equation}
\label{eq:app-stage001-a-moment}
a(t)=\frac{1}{4\pi}\int_{S^2}R(\Omega,0,t)\,d\Omega,
\end{equation}
and, if \(W_b(\Omega,t)\) is defined implicitly by
\begin{equation}
R(\Omega,W_b(\Omega,t),t)=0,
\end{equation}
then
\begin{equation}
\label{eq:app-stage001-l-moment}
L(t)=\frac{1}{4\pi}\int_{S^2}W_b(\Omega,t)\,d\Omega.
\end{equation}
Thus the old closure survives as a moment projection of the distributed wall field.

\paragraph{3. Promoted confinement coupling.}
The old confinement dependence \(V_{\rm conf}(\mathbf X;a,L)\) is promoted to a moving-surface dependence
\begin{equation}
\label{eq:app-stage001-promoted-conf}
V_{\rm conf}(\mathbf X;\Sigma)=V_{\rm wall}\!\left(\frac{\Sigma(\mathbf X,t)}{\ell_c}\right),
\end{equation}
where \(\ell_c\) is the wall thickness.  Around the reference branch,
\begin{equation}
V_0(\mathbf X)=V_{\rm wall}\!\left(\frac{\Sigma_0(\mathbf X)}{\ell_c}\right).
\end{equation}
The first variation is
\begin{equation}
\label{eq:app-stage001-delta-v-general}
\delta V_{\rm conf}
=\frac{V_{\rm wall}'(\Sigma_0/\ell_c)}{\ell_c}\,\delta\Sigma.
\end{equation}
Since \(\Sigma=r-R\), a wall perturbation \(R=R_0+\eta\) gives
\begin{equation}
\label{eq:app-stage001-delta-sigma}
\delta\Sigma=-\eta(\Omega,w,t),
\end{equation}
and therefore
\begin{equation}
\label{eq:app-stage001-delta-v}
\boxed{
\delta V_{\rm conf}
=-\frac{V_{\rm wall}'(\Sigma_0/\ell_c)}{\ell_c}\,\eta(\Omega,w,t)}.
\end{equation}
This is the first direct wall-to-BdG coupling used later.

\paragraph{4. Harmonic decomposition.}
Write
\begin{equation}
\label{eq:app-stage001-wall-perturbation}
R(\Omega,w,t)=R_0(w)+\eta(\Omega,w,t).
\end{equation}
Expand \(\eta\) in real spherical harmonics:
\begin{equation}
\label{eq:app-stage001-harmonic-expansion}
\eta(\Omega,w,t)
=\eta_0(w,t)Y_{00}(\Omega)
+\sum_{m\in P_2({\rm real})}q_{2m}(w,t)Y^{\rm real}_{2m}(\Omega)
+\eta_{\ge3}(\Omega,w,t).
\end{equation}
The real \(P_2\) set is
\begin{equation}
\label{eq:app-stage001-p2-set}
P_2({\rm real})=\{20,21c,21s,22c,22s\}.
\end{equation}
With the normalized monopole
\begin{equation}
Y_{00}=\frac{1}{2\sqrt\pi},
\qquad
\int_{S^2}Y_{00}^2\,d\Omega=1,
\end{equation}
we have
\begin{equation}
\frac{1}{4\pi}\int_{S^2}Y_{00}\,d\Omega=\frac{1}{2\sqrt\pi}.
\end{equation}
Therefore the normalized monopole coefficient and the physical mouth-average shift are related by
\begin{equation}
\label{eq:app-stage001-q00-da}
q_{00}(0,t)=2\sqrt\pi\,\delta a(t).
\end{equation}
A useful axisymmetric split is
\begin{equation}
\label{eq:app-stage001-axisymmetric-split}
\eta_0(w,t)=\alpha_a(w)\delta a(t)+\alpha_L(w)\delta L(t)+g(w,t),
\end{equation}
where \(g\) is the residual axisymmetric geometry lane orthogonal to the \((a,L)\) collective moments.

\paragraph{5. Minimal distributed wall action.}
The working passive quadratic wall action is
\begin{equation}
\label{eq:app-stage001-wall-action}
S_{\eta}^{(2)}
=\frac12\int dt\,dw\,d\Omega\,\sqrt{\gamma_0}
\left[
\mu_\eta(w)(\partial_t\eta)^2
-T_w(w)(\partial_w\eta)^2
-T_\Omega(w)\eta(-\Delta_{S^2})\eta
-K_\eta(w)\eta^2
\right].
\end{equation}
The coefficients \(\mu_\eta,T_w,T_\Omega,K_\eta\) are branch constitutive functions.  From this point onward the ledger adopts the densitized one-dimensional convention: the background surface weight \(\sqrt{\gamma_0}\) is absorbed into the effective axial coefficients and mode profiles, so the reduced formulas are written with respect to \(dw\).  The weighted surface form is still legitimate, but then the Euler--Lagrange operator carries the extra first-derivative term coming from \(\partial_w\sqrt{\gamma_0}\); the SymPy and Mathematica audits check both forms explicitly.

Projecting onto a harmonic \(Y_{lm}\), using
\begin{equation}
-\Delta_{S^2}Y_{lm}=l(l+1)Y_{lm},
\end{equation}
gives the modal wall operator
\begin{equation}
\label{eq:app-stage001-modal-wall-pde}
\boxed{
\mu_\eta\,\partial_t^2q_{lm}
-\partial_w\!\bigl(T_w\partial_wq_{lm}\bigr)
+\bigl[K_\eta+l(l+1)T_\Omega\bigr]q_{lm}
=S_{lm}^{(\psi,A)}+f_{lm}^{\rm ext}}.
\end{equation}
The scalar lane \(l=0\) and the grouped real quadrupole lane \(l=2\) are therefore separated before any additional closure.

\paragraph{6. Linearized matter and gauge sectors.}
Choose stationary background fields
\begin{equation}
\psi=e^{-i\mu_0t/\hbar}\psi_0(\mathbf X),
\qquad
A_M=A_{M0}(\mathbf X),
\qquad
R=R_0(w).
\end{equation}
Let
\begin{equation}
\psi=e^{-i\mu_0t/\hbar}\bigl(\psi_0+\delta\psi\bigr),
\qquad
A_M=A_{M0}+\delta A_M.
\end{equation}
The matter perturbation satisfies a BdG-type system
\begin{equation}
\label{eq:app-stage001-bdg-schematic}
i\hbar\partial_t
\begin{pmatrix}\delta\psi\\ \delta\psi^*\end{pmatrix}
=\mathcal L_{\rm BdG}
\begin{pmatrix}\delta\psi\\ \delta\psi^*\end{pmatrix}
+\mathcal C_A[\delta A_M]
+\mathcal C_\eta[\eta],
\end{equation}
where \(\mathcal C_\eta\) contains the coupling \eqref{eq:app-stage001-delta-v}.  The localized Maxwell equation linearizes to
\begin{equation}
\label{eq:app-stage001-linear-maxwell}
\boxed{
\partial_M\bigl(Z(w)\delta F^{MN}\bigr)
+\frac1\xi\partial^N(\partial\!\cdot\!\delta A)
=\mu_0\delta J^N}.
\end{equation}
At this stage the mixed variables
\begin{equation}
\delta A_w,
\qquad
\delta J^w,
\qquad
\delta F_{\mu w}
\end{equation}
are retained.  Their later suppression in the far-field brane limit is a reduction, not an ontology deletion.

\paragraph{7. Mouth/worldtube response operator.}
At the mouth \(w=0\), choose the real port basis
\begin{equation}
P_{00}=Y_{00},
\qquad
P_{20},P_{21c},P_{21s},P_{22c},P_{22s}.
\end{equation}
Let \(u_B(\omega)\) denote projected effort variables and \(j_A(\omega)\) the corresponding generalized fluxes or tractions.  The reduced response operator is written
\begin{equation}
\label{eq:app-stage001-response-operator}
j_A(\omega)=\sum_B Z_{{\rm eff},AB}(\omega)u_B(\omega).
\end{equation}
On an isotropic reference throat, this operator block-diagonalizes into scalar/geometry, grouped real \(P_2\), and higher-\(l\) sectors.

\stagefield{Output}{The stage outputs the shape-field closure \eqref{eq:app-stage001-sigma}, the confinement variation \eqref{eq:app-stage001-delta-v}, the harmonic split \eqref{eq:app-stage001-harmonic-expansion}, the modal wall PDE \eqref{eq:app-stage001-modal-wall-pde}, and the response-operator target \eqref{eq:app-stage001-response-operator}.  These are the definitions without which later grouped response coefficients have no PDE-side interpretation.}

\stagefield{Checks}{\begin{verificationchecklist}
\item The moment check follows from the normalized monopole integral and gives \(q_{00}=2\sqrt\pi\,\delta a\).
\item The sign check for \(\delta V_{\rm conf}\) follows from \(\delta\Sigma=-\eta\).
\item The harmonic selection check follows from orthogonality on \(S^2\), so the isotropic linear problem has no \(l=0\leftrightarrow l=2\) mixing.
\item The ontology check keeps \(A_w,J^w,F_{\mu w}\) alive before any far-field brane reduction.
\end{verificationchecklist}}

\stagefield{Downstream use}{Stage~002 uses the same wall action to recover \((a,L)\).  Stage~003 uses \eqref{eq:app-stage001-delta-v} as the wall--BdG coupling.  Stages~004--021 use the retained mixed variables as the port-active gauge sector and reduce them to the one-port normal form.  Stages~022--023 use the grouped real \(P_2\) response operator to formulate the outgoing-normalization product.}
